
\documentclass{article}
\usepackage{colortbl}
\usepackage{makecell}
\usepackage{multirow}
\usepackage{supertabular}

\begin{document}

\newcounter{utterance}

\centering \large Interaction Transcript for game `cladder', experiment `full\_v1.5\_default', episode 4490 with Qwen3.5{-}9B.
\vspace{24pt}

{ \footnotesize  \setcounter{utterance}{1}
\setlength{\tabcolsep}{0pt}
\begin{supertabular}{c@{$\;$}|p{.15\linewidth}@{}p{.15\linewidth}p{.15\linewidth}p{.15\linewidth}p{.15\linewidth}p{.15\linewidth}}
    \# & \multicolumn{2}{c}{Player} && \multicolumn{2}{c}{Game Master} \\
    \hline

    \theutterance \stepcounter{utterance}  
    & & & \multicolumn{4}{p{0.6\linewidth}}{
        \cellcolor[rgb]{0.9,0.9,0.9}{
            \makecell[{{p{\linewidth}}}]{
                \texttt{\tiny{[P1$\langle$GM]}}
                \texttt{You are an expert in causal inference. The following question is not a typical commonsense query, but rather a meticulously designed question created by a professor specializing in causal inference, intended to assess the students' mastery of the course content.} \\
\texttt{Imagine a self{-}contained, hypothetical world with only the following conditions, and without any unmentioned factors or causal relationships: Listening to jazz has a direct effect on tar deposit. Tar deposit has a direct effect on lung cancer. For people who do not listen to jazz, the probability of lung cancer is 61\%. For people who listen to jazz, the probability of lung cancer is 79\%. For people who listen to jazz, would it be less likely to see lung cancer if the person did not listen to jazz?} \\
\texttt{Start your answer with 'Yes' or 'No', followed by additional reasoning or evidence to support your explanation.} \\
            }
        }
    }
     \\ \\

    \theutterance \stepcounter{utterance}  
    & \multicolumn{4}{p{0.6\linewidth}}{
        \cellcolor[rgb]{0.561,0.737,0.561}{
            \makecell[{{p{\linewidth}}}]{
                \texttt{\tiny{[P1$\rangle$GM]}}
                \texttt{No} \\
\\ 
\texttt{**Reasoning:**} \\
\\ 
\texttt{To determine the answer, we must analyze the causal structure and the specific probabilities provided within the hypothetical world.} \\
\\ 
\texttt{1.  **Causal Structure**: The problem defines a directed acyclic graph (DAG) where:} \\
\texttt{    *   \$J\$ = Listening to Jazz} \\
\texttt{    *   \$T\$ = Tar Deposit} \\
\texttt{    *   \$L\$ = Lung Cancer} \\
\texttt{    *   The paths are \$J \textbackslash{}rightarrow T \textbackslash{}rightarrow L\$.} \\
\texttt{    *   Crucially, the prompt states there are "no unmentioned factors or causal relationships." This implies there is no confounding variable (like smoking) that causes both Jazz listening and Tar deposits. Therefore, the association between Jazz and Lung Cancer is entirely mediated by Tar deposits.} \\
\\ 
\texttt{2.  **Data Analysis**:} \\
\texttt{    *   \$P(L | \textbackslash{}neg J) = 61\textbackslash{}\%\$ (Probability of lung cancer for non{-}listeners).} \\
\texttt{    *   \$P(L | J) = 79\textbackslash{}\%\$ (Probability of lung cancer for listeners).} \\
\\ 
\texttt{3.  **Counterfactual Evaluation**:} \\
\texttt{    *   The question asks: "For people who listen to jazz (\$J\$), would it be less likely to see lung cancer if the person did not listen to jazz (\$\textbackslash{}neg J\$)?"} \\
\texttt{    *   Mathematically, this asks if \$P(L | J, \textbackslash{}neg J) < P(L | J)\$.} \\
            }
        }
    }
    & & \\ \\

    \theutterance \stepcounter{utterance}  
    & & & \multicolumn{2}{p{0.3\linewidth}}{
        \cellcolor[rgb]{0.9,0.9,0.9}{
            \makecell[{{p{\linewidth}}}]{
                \texttt{\tiny{[GM$|$GM]}}
                \texttt{No} \\
            }
        }
    }
    & & \\ \\

    \theutterance \stepcounter{utterance}  
    & & & \multicolumn{2}{p{0.3\linewidth}}{
        \cellcolor[rgb]{0.9,0.9,0.9}{
            \makecell[{{p{\linewidth}}}]{
                \texttt{\tiny{[GM$|$GM]}}
                \texttt{yes} \\
            }
        }
    }
    & & \\ \\

    \theutterance \stepcounter{utterance}  
    & & & \multicolumn{2}{p{0.3\linewidth}}{
        \cellcolor[rgb]{0.9,0.9,0.9}{
            \makecell[{{p{\linewidth}}}]{
                \texttt{\tiny{[GM$|$GM]}}
                \texttt{game\_result = LOSE} \\
            }
        }
    }
    & & \\ \\

\end{supertabular}
}

\end{document}
